\documentclass[acmsmall,11pt,screen,nonacm]{acmart}
\geometry{letterpaper,margin=1in,includehead=false,includefoot=false,footskip=24pt}
\settopmatter{printacmref=false,printccs=false,printfolios=true}
\setcopyright{none}
\renewcommand{\bibfont}{\normalfont\normalsize}
\newcommand{\dist}{\operatorname{dist}}
\newcommand{\Ot}{\widetilde{O}}
\newcommand{\F}{\mathbb{F}}
\newcommand{\cP}{\mathcal{P}}
\newtheorem{theorem}{Theorem}[section]
\newtheorem{lemma}[theorem]{Lemma}
\newtheorem{proposition}[theorem]{Proposition}
\title{Polynomial-Time Construction of Vertex-Fault-Tolerant Mixed-Stretch Spanners}
\author{Zijian Zeng}
\affiliation{\institution{UCSI University}\country{Malaysia}}
\email{1002266693@ucsiuniversity.edu.my}
\date{}
\begin{document}
\begin{abstract}
\normalsize
Parter, Popova and Tzalik construct vertex-fault-tolerant spanners with
$(O_\delta(\sqrt{k}),O_\delta(k))$ stretch and
$\widetilde{O}_\delta(f^{1+\delta}n^{1+1/k})$ edges, for every fixed
$\delta\in(0,1)$. Their construction at this tradeoff enumerates short
paths and has running time about $fn^{\sqrt{k}}$.
We give a deterministic polynomial-time construction at the same tradeoff,
uniformly over the fault and stretch parameters. More precisely, the edge
bound is $\widetilde{O}_\delta((f+1)^{1+\delta}n^{1+1/k})$, including $f=0$;
the polynomial exponent depends only on $\delta$.

The algorithm uses established representative-family machinery to retain
short paths against every fault set, and a protected-path version of the
length-bounded-cut test of Dinitz and Robelle. This computation costs
$n^{O(1)}(f+d+1)^{O(d+1)}$ for path length $d$.
After capping $k$ at $O(\log n)$, it is polynomial when
$f^{\eta d}<n$, where $d=\lfloor\sqrt{k}\rfloor$ and
$\eta=\min\{\delta/2,1/4\}$. In the complementary regime, the known
multiplicative construction of Dinitz and Robelle already meets the desired
density bound. We verify that the insertion-witness analysis of Parter,
Popova and Tzalik applies to the resulting partial path scan and track its
parameter factors explicitly. Pendant-path reductions then give the
guarantee simultaneously for every surviving pair and fault set.
The resulting algorithm realizes the full mixed-stretch tradeoff in
deterministic polynomial time.
\end{abstract}
\maketitle
\clearpage

\section{The construction gap and its resolution}

Let $G=(V,E)$ be an undirected unweighted simple graph on $n$ vertices.
For $F\subseteq V$, the graph $G-F$ deletes $F$ and its incident edges.
A subgraph $H\subseteq G$ is an $f$-vertex-fault-tolerant
$(\alpha,\beta)$-spanner if, for every $|F|\le f$ and $u,v\notin F$ with
finite $\dist_{G-F}(u,v)$,
\[
 \dist_{H-F}(u,v)\le \alpha\dist_{G-F}(u,v)+\beta.
\]
Thus the fault set is chosen after the subgraph, and one subgraph must
handle all faults simultaneously.

Parter, Popova and Tzalik~\cite{PPT26} establish an almost-linear dependence
on $f$ for sublinear multiplicative stretch: their Theorem~1.4 gives
$\Ot_\delta(f^{1+\delta}n^{1+1/k})$ edges and
$(O_\delta(\sqrt{k}),O_\delta(k))$ stretch. Their introduction identifies
improving this construction's running time as a remaining problem. Its
short-path enumeration takes about $fn^{\sqrt{k}}$ time. Their polynomial
construction with a larger polynomial dependence on $f$ has a different
tradeoff. Here we make the almost-linear-fault tradeoff polynomial-time
constructible, without requiring either $f$ or $k$ to be fixed.

\begin{theorem}\label{thm:main}
For every fixed rational $\delta\in(0,1)$ there are constants
$C_\delta,b_\delta>0$ and a deterministic algorithm, polynomial in the
explicit input length, that takes $G$ and binary integers $f\ge0,k\ge1$
and returns an $f$-vertex-fault-tolerant
$(C_\delta\sqrt{k},C_\delta k)$-spanner with at most
\[
 C_\delta(f+1)^{1+\delta}n^{1+1/k}\log(n+2)^{b_\delta}
\]
edges. The running-time exponent depends only on $\delta$.
\end{theorem}

Every fixed real $\delta\in(0,1)$ follows by choosing a smaller positive
rational exponent. Throughout, $\Ot_\delta$ suppresses only a bounded
power of $\log(n+2)$ and constants depending on $\delta$; it does not
suppress a factor exponential in $\sqrt{k}$.

\paragraph{Why complementary regimes help.}
The algorithm's more expensive branch computes fault-representative paths
of length $d=\lfloor\sqrt{k}\rfloor$. Standard exterior-algebra reduction
gives a cost of $n^{O(1)}(f+d+1)^{O(d+1)}$; it does not by itself give a
uniform polynomial-time algorithm. Set $\eta=\min\{\delta/2,1/4\}$ and
cap $k$ at $O(\log n)$. If $f^{\eta d}<n$, this cost becomes polynomial
with exponent depending only on $\eta$. If $f^{\eta d}\ge n$, a known
$(2d-1)$-multiplicative spanner with
$O(df^{1-1/d}n^{1+1/d})$ edges already has the required density. The
larger allowed fault factor pays for its larger vertex-density exponent.
This switch, together with a proof that the actual partial scan retains
the source packing bound, supplies the complete construction.

\paragraph{Prior tools and the precise delta.}
Dinitz and Robelle~\cite{DR20} give a deterministic multiplicative
construction and an elementary length-bounded-cut approximation. We use
their Theorems~2 and~9 for size and time, and adapt their cut test by
protecting all vertices of a candidate path. Bodwin, Dinitz and
Robelle~\cite{BDR21} give optimal-size polynomial-time multiplicative
spanners; the simpler construction suffices here because a factor
polynomial in $k$ is allowed. Neither multiplicative theorem alone supplies
the mixed-stretch tradeoff in every regime.

Fomin, Lokshtanov, Panolan and Saurabh~\cite{FLPS16} develop the
representative-family algebra and its path dynamic-programming
applications. We give the elementary uniform-matroid specialization needed
here. Fault-avoiding short-path coverings are also established machinery:
Karthik C.~S. and Parter~\cite{KP23} include a vertex-fault variant;
the more recent replacement-path-covering results
of~\cite{BCCCS26,BCCS26} improve covering bounds in edge-fault settings.
These works supply the established short-path-covering tools. The sparsity
analysis relies on the path systems, spread families and robust-sunflower
counting of Parter, Popova and Tzalik~\cite{PPT26}. The additional work
here is their use for the partial scan below, with uniform parameter
accounting, and the complementary-regime construction theorem.

\section{Exact fault-representative paths}\label{sec:represent}

Fix an $N$-vertex graph $J$, a length $d<N$, and $f\ge0$. We compute a
family for each pair $x,y$ with the following property: whenever a simple
length-$d$ $x$--$y$ path avoids a fault set of size at most $f$, some
retained length-$d$ $x$--$y$ path also avoids that set.

Put $T=\min\{N,f+d+1\}$. Choose a prime $p\in[N,2N)$ by trial division
and distinct elements $\alpha_v\in\F_p$ for the vertices. This takes
polynomial time in $N$. Let $a_v=(1,\alpha_v,\ldots,\alpha_v^{T-1})^\top$.
Any at most $T$ distinct columns are independent by the Vandermonde
determinant. For an $h$-set $A$, let $w(A)$ be the exterior product of its
columns, in vertex-label order, represented by all $h$-row minors. Its
dimension is $\binom{T}{h}$.

\begin{lemma}[Standard representative reduction]\label{lem:basis}
For a family $\mathcal A$ of $h$-sets, retain actual members whose vectors
form a basis of $\operatorname{span}\{w(A):A\in\mathcal A\}$. The retained
family has size at most $\binom{T}{h}$. For every $Y$ of size at most $T-h$,
if some member of $\mathcal A$ avoids $Y$, a retained member avoids $Y$.
\end{lemma}
\begin{proof}
If $A\cap Y=\varnothing$, independence of the at-most-$T$ columns gives
$w(A)\wedge w(Y)\ne0$. Expressing $w(A)$ in the retained basis shows that
at least one retained vector has nonzero product with $w(Y)$. Its actual
support avoids $Y$, since a repeated column would give zero. This is an
exact finite-field argument, a specialization of~\cite{FLPS16}.
\end{proof}

For each start $x$, store the path $(x)$ at $x$, with empty families at
other endpoints. At each stage append a neighbor to every stored path,
discard repetitions, and apply Lemma~\ref{lem:basis} separately at each
new endpoint. Always keep actual paths, not only their vectors.

\begin{lemma}\label{lem:paths}
After stage $i$, the endpoint family represents all simple $i$-edge paths
from $x$ against sets of size at most $T-i-1$. In particular the final
families have the stated fault-avoidance property. All families, for all
starts, can be computed in
\[
 N^{O(1)}(f+d+1)^{O(d+1)}
\]
bit operations, with absolute constants in the $O$ exponents.
\end{lemma}
\begin{proof}
The invariant holds at stage zero. Suppose a path extended to $v$ avoids
$Y$, where $|Y|\le T-i-2$. Apply the stage-$i$ invariant to
$Y\cup\{v\}$: a stored parent avoids this set, can be extended to $v$
without repetition, and avoids $Y$. Basis reduction preserves that
property. This proves the induction.

At the final stage, if $T=f+d+1$, the avoidance budget is $f$. If $T=N$,
any fault set disjoint from a $(d+1)$-vertex path has size at most
$N-d-1=T-d-1$, so this rank cap loses no necessary query. In an undirected
graph it suffices to retain the orientation from the smaller endpoint.

Let $L=\max_{1\le h\le d+1}\binom{T}{h}\le(f+d+1)^{d+1}$. There are at
most $L$ stored paths per endpoint. A stage has at most $N^2L$ candidate
extensions per start. Computing minors by elimination and reducing their
vectors takes a fixed polynomial in $N,d,L,\log N$; for example
$O(N^3(d+1)^4L^3)$ field operations suffices for all starts and stages.
Field elements have $O(\log N)$ bits. This proves the bound without
enumerating all simple paths or fault sets.
\end{proof}

\section{A protected-path scan}\label{sec:scan}

Fix an integer $r\ge d$ and start with the empty subgraph $H$ of $J$.
Scan each retained path $P$ once, in a fixed computable order. With endpoints
$x,y$ and protected set $S=V(P)$, perform the following test.
Initialize $C=\varnothing$. For $j=0,\ldots,f$, run BFS in $H-C$.
If its shortest $x$--$y$ path has length greater than $r$ (including
infinite distance), return $C$ and insert all edges of $P$.
Otherwise choose such a shortest path $Q$ and set $B=V(Q)\setminus S$.
If $B=\varnothing$ or $j=f$, reject $P$. In every other case replace
$C$ by $C\cup B$ and continue.

\begin{lemma}\label{lem:scan}
Every insertion has a witness $C$ disjoint from $P$, with
$|C|\le(r-1)f$ and $\dist_{H-C}(x,y)>r$ immediately before insertion.
Every rejection certifies that
\[
 \dist_{H-F}(x,y)\le r
 \quad\text{for all } |F|\le f,\ F\cap V(P)=\varnothing.
\]
The final graph satisfies this distance bound for every pair admitting a
length-$d$ path in $J-F$.
\end{lemma}
\begin{proof}
All collected vertices lie outside $S$. Each short $Q$ has at most $r-1$
unprotected vertices, since its two endpoints lie in $S$. Before any return
at most $f$ sets were collected. The final BFS supplies the witness.
An insertion adds a new edge: otherwise $P$ itself would survive $C$ in
the previous graph and have length $d\le r$.

If $B$ is empty, $Q$ survives every fault set disjoint from $S$. Otherwise
a rejection on the last round produces $f+1$ nonempty sets $B$. They are
pairwise disjoint because later BFS paths avoid the earlier collected
sets. Any fault set outside $S$ destroying every short path must intersect
each $B$, impossible with at most $f$ vertices. This proves the rejection
certificate, the protected version of the cut argument in~\cite{DR20}.

Now fix $F$ and a length-$d$ path avoiding it. Lemma~\ref{lem:paths}
gives a retained path avoiding $F$. If inserted it survives in the output;
if rejected, the certificate gives a short route in the graph at that
time. Edge insertions preserve the route. Since this argument fixes
arbitrary $F$, it proves the simultaneous guarantee.
\end{proof}

The tests cost at most $O(N^2L(f+1)(N+|E(J)|))$ in addition to the
representative computation. Denote the inserted paths by $P_1,\ldots,P_m$
and the witnesses by $C_1,\ldots,C_m$. With $H_i=\bigcup_{j<i}P_j$,
\begin{equation}\label{eq:witness}
 C_i\cap V(P_i)=\varnothing,\qquad |C_i|\le(r-1)f,\qquad
 \dist_{H_i-C_i}(\operatorname{ends}(P_i))>r.
\end{equation}
We next count this actual sequence. There is no containment assertion
between its output and a different complete greedy run.

\section{Transferring the insertion-witness count}\label{sec:packing}

The following consequence uses the deep counting machinery of~\cite{PPT26}.
We supply the transfer and parameter accounting, and identify the precise
prior results used; we do not reprove their park or robust-sunflower theorems.

\begin{proposition}\label{prop:packing}
Fix $0<\eta\le1/4$, $k\ge4$, $d=\lfloor\sqrt{k}\rfloor$ and
$r=\lceil12k/\eta\rceil$, with $k=O(\log(N+2))$ and an absolute constant
in this last bound. For any sequence of simple $d$-edge paths with
disjoint insertion witnesses as in~\eqref{eq:witness}, but with
$|C_i|\le q$ for any integer $q\ge1$, their union has
\[
 \Ot_\eta(q^{1+\eta}N^{1+1/k})
\]
edges. The logarithmic exponent is bounded independently of $q,k,N$.
\end{proposition}
\begin{proof}
If $q\ge N$, simplicity gives $|E(H)|\le N^2\le qN$, so suppose $q<N$.
Put
\[
 a=\sqrt{k},\quad t=\lceil a\rceil,\quad
 R=\lceil a/\eta\rceil,\quad s=2R-1,\quad
 D=\left\lceil12C's^3q^{1+\eta}N^{1/k}\right\rceil,
\]
where $C'\ge2$ is the fixed absolute constant required in
\cite[Lemma 5.14]{PPT26}. In particular $D\ge4s^2q$.
The rounding bounds needed below are
\begin{equation}\label{eq:rounding}
 ds+2R\le4k/\eta,\qquad 2R+d\le3a/\eta,\qquad
 (2t-1)(2R+d)\le8k/\eta.
\end{equation}
They follow from $R\le a/\eta+1$, $d\le a$, $t\le a+1$, $a\ge2$ and
$\eta\le1/4$. Thus $ds\le r$.

Call an edge $uv$ on $P_i$ $s$-critical if
$\dist_{H_i-C_i}(u,v)>s$. Such an edge exists: otherwise concatenate
the $d$ detours to contradict~\eqref{eq:witness}. This is exactly the
argument of~\cite[Claim 4.3]{PPT26}. In particular this edge is absent
from $H_i$ and both its endpoints avoid $C_i$.

\paragraph{Why the analytical path system still applies.}
Run Algorithm~2 of~\cite{PPT26} analytically, with parameter $q$, on these
critical edges and witnesses. It is not a step of our implementation.
The induction in its Lemma~4.12 takes the existing path system, a critical
edge, a compatible fault set, and the prepend/filter update. Its invariants
I1--I7 assert containment in $H_i$, collective bounds on paths through
vertex sets, the pairwise park property, bounded total size, progress per
boost, and a bounded number of boosts. The path-exchange Lemma~4.16 is for
an arbitrary graph and critical pair. These statements and their induction
use no rejected or omitted paths, no completeness of an endpoint scan,
and no shortestness assumption on the inserted path. The critical edge's
absence supplies distinct first-neighbor insertions exactly as in that
induction. Consequently its local/global partition applies to our sequence.
By its Theorem~4.5, the number of local rounds is
\begin{equation}\label{eq:local}
 O(NDs)=O_\eta(s^4q^{1+\eta}N^{1+1/k}).
\end{equation}

\paragraph{Centers with uniform factors.}
Proposition~4.14 and Appendix~A of~\cite{PPT26} freeze the first full
$R$-path family at each vertex and obtain a center set $S$. Every global
round has two nonempty families of length-$R$ paths in $H_i$ from middle
vertices of $P_i$ to $S$. Each is $C'sq$-spread: for a family
$\mathcal A$ based at $v$ and every $U\subseteq V\setminus\{v\}$,
\[
 |\{Q\in\mathcal A:U\subseteq V(Q)\}|\le
 |\mathcal A|(C'sq)^{-|U|}.
\]
The same proposition bounds the prefix and suffix distances along $P_i$
in $H_i-C_i$ by their respective lengths times $s$.
Its proof uses the analytical path system and the present witnesses, so
the transfer just verified supplies its hypotheses.

It is necessary to keep the constants in its sampling bound. With
\[
 M=\log(100\cdot2^{R+2}D^RN),\qquad
 p=\min\{11(10C'sq/D)^RM,1\},
\]
Lemma~A.1 and the proof of Proposition~4.14 give
\begin{equation}\label{eq:centers}
 |S|\le22M q^{-\eta R}N^{1-R/k}.
\end{equation}
Indeed,
\[
 10C'sq/D\le (5/(6s^2))q^{-\eta}N^{-1/k}.
\]
The factor $(5/(6s^2))^R$ can be dropped because it is at most one.
If the probability is capped at one, the uncapped expression is at least
one, so the same bound permits taking all $N$ vertices as centers.
Since $q<N$, $\log D=O_\eta(\log(N+2))$ and
$M=O_\eta(\log(N+2)^{3/2})$. Thus no positive factor exponential in
$\sqrt{k}$ has been absorbed into logarithmic notation. The sampled
centers are existential witnesses for counting, not randomness in the
construction algorithm.

\paragraph{Global rounds.}
Observation~5.8 and Lemma~5.9 of~\cite{PPT26} convert these families into
an evolving virtual multigraph: each virtual edge has a realizing path
in $H_{i+1}$ avoiding $C_i$, and its prior virtual distance under $C_i$
exceeds $2t-1$. Their distance contradictions remain valid by
\eqref{eq:rounding} and $r\ge12k/\eta$. Lemma~5.14 converts the spread
families into the robust sunflower whose kernel is the middle subpath
of $P_i$. Its size is at most $w=d+1$. This step uses the spreadness and
compatibility conditions already stated. Theorem~5.15, which counts
arbitrary iterative robust cluster-fault-tolerant inputs, then bounds
the number of global rounds by
\[
 O(wq^w|S|^{1+1/t}).
\]
Substituting~\eqref{eq:centers}, the exponents satisfy
\[
 w-\eta R(1+1/t)\le1,\qquad
 (1-R/k)(1+1/t)\le1.
\]
The first uses $\eta R\ge a\ge d$, and the second uses
$R/k\ge1/(t+1)$. Also $1+1/t\le3/2$, so the positive factor
$(22M)^{1+1/t}$ remains polylogarithmic. Hence the global count is
$\Ot_\eta(wqN)$. If $S$ is empty there are no global rounds, since their
families must be nonempty. Combining this count with~\eqref{eq:local},
and charging at most $d$ edges per round, proves the proposition. All
remaining factors in $s,d,w$ are polynomial in $k=O(\log(N+2))$.
\end{proof}

For the scan, take $q=(r-1)f$ when $f\ge1$, and $q=1$ when $f=0$.
Empty witnesses are admissible with the larger analytical budget. Thus
its output has
\begin{equation}\label{eq:dtor-size}
 \Ot_\eta\bigl((1+(r-1)f)^{1+\eta}N^{1+1/k}\bigr)
\end{equation}
edges. The difficult combinatorial ingredients of this bound remain
the cited results of Parter, Popova and Tzalik; the transfer checks why
their hypotheses hold for the construction actually used here.

\section{All distances and simultaneous vertex faults}\label{sec:scales}

For $\ell=0,\ldots,d-1$, form $J_\ell$ by attaching a private
$\ell$-edge pendant chain to every vertex of the base graph $G$.
It has $N_\ell=(\ell+1)n\le dn$ vertices. Run the scan on each augmented
graph with the same $d,r,f$, retain its edges with both endpoints in
$V(G)$, and let $H$ be the union of these projections. This is the
pendant reduction of~\cite[Appendix B]{PPT26}; we spell out its fault
quantifiers.

Fix $F\subseteq V(G)$ with $|F|\le f$. If $u,v$ have surviving base
distance $b=d-\ell\ge1$, the pendant tip $v'$ at $v$ in $J_\ell-F$ has
distance exactly $b+\ell=d$ from $u$. Pendant chains give no shortcuts.
Lemma~\ref{lem:scan} supplies a route of length at most $r$ in the augmented
output minus this same $F$. A shortest such route is simple. It must reach
$v$ and traverse the private chain, and cannot make excursions into other
pendant chains. Its base segment is therefore a route in the projection
minus $F$, of length at most $r-\ell$. The augmented fault guarantee
includes $F$ itself; no additional or substituted faults are needed.

Segment a shortest $G-F$ path into $d$-edge pieces and a final remainder
of length $b<d$. Every piece is geodesic in $G-F$. The $\ell=0$ projection
routes full pieces within $r$; for $b>0$ use $\ell=d-b$ to route the
remainder within $r-(d-b)\le r$. Thus, for all surviving pairs,
\begin{equation}\label{eq:stretch}
 \dist_{H-F}(u,v)\le(r/d)\dist_{G-F}(u,v)+r.
\end{equation}
By~\eqref{eq:dtor-size} the total number of edges is at most
\begin{equation}\label{eq:scales-size}
 d\,\Ot_\eta\bigl((1+(r-1)f)^{1+\eta}(dn)^{1+1/k}\bigr).
\end{equation}
Only base edges remain, so this is an unweighted subgraph spanner.

\section{Uniform running time and boundary cases}\label{sec:switch}

\begin{proof}[Proof of Theorem~\ref{thm:main}]
For $n\le2$, return $G$. Otherwise cap $f$ at $n-2$. This loses no
requirement for distinct surviving pairs: a fault set with two surviving
vertices already has size at most $n-2$. Cap the input $k$ at
$K=\lceil\log_2(n+2)\rceil$. For an original $k>K$, a construction at
$K$ changes the density factor by at most two since $n^{1/K}\le2$;
its stretch terms only improve. Below, $k$ denotes the capped parameter.
Reading and capping arbitrarily large binary integers is polynomial in
the explicit input length.

Set $\eta=\min\{\delta/2,1/4\}$. We use the following deterministic
theorem of Dinitz and Robelle~\cite[Theorems 2 and 9]{DR20}: for
integers $t\ge1$ and $g\ge1$ there is a $g$-VFT $(2t-1)$-spanner with
\[
 O(tg^{1-1/t}n^{1+1/t})\ \text{edges},\qquad
 O(mt g^{2-1/t}n^{1+1/t})\ \text{construction time}.
\]
For $k\le3$, use it with $t=k$ and $g=\max\{1,f\}$.
The density meets the requested bound, $2k-1\le3\sqrt{k}$ in this
range, and the runtime is polynomial. In particular $f=0$ does not
justify returning a dense graph in this branch.

Suppose $k\ge4$, and put $d=\lfloor\sqrt{k}\rfloor$,
$r=\lceil12k/\eta\rceil$. If $f\ge1$ and
\begin{equation}\label{eq:switch}
 f^{\eta d}\ge n,
\end{equation}
use the multiplicative construction with $t=d,g=f$. From
$n^{1/d}\le f^\eta$, its edge count is
\[
 O(df^{1-1/d}n^{1+1/d})
 \le O(df^{1+\eta-1/d}n)
 \le O(d(f+1)^{1+\delta}n^{1+1/k}).
\]
The factor $d\le\sqrt K$ is logarithmic slack. Its stretch is
$2d-1\le2\sqrt{k}$, and its time is polynomial since $f\le n$,
$d\le\sqrt K$ and $m\le n^2$. The switch is exact: if
$\eta=A/B$ in lowest terms, compare the integers $f^{Ad}$ and $n^B$.

Otherwise use the pendant construction of Section~\ref{sec:scales}.
Here $f=0$ or $\log f<\log n/(\eta d)$. Since
$f+d+1\le(f+1)(d+1)$,
\begin{align*}
 (d+1)\log(f+d+1)
 &\le(d+1)\log(f+1)+(d+1)\log(d+1)\\
 &=O_\eta(\log(n+2)).
\end{align*}
For $f\ge1$ this uses $\log(f+1)\le\log2+\log f$ and
$(d+1)/d\le2$. For $f=0$ the first term vanishes. In both cases
$(d+1)\log(d+1)\le2d^2\le2K$, with natural logarithms.
Lemma~\ref{lem:paths} and the scan cost are consequently polynomial in
$n$, with exponent depending only on $\eta$. There are at most
$d\le\sqrt K$ augmented graphs, each with $N\le dn\le n^2$;
the switch uses the base $n$. BFS introduces only polynomial factors.
In particular no fault enumeration occurs in the algorithm.

In~\eqref{eq:scales-size}, all positive factors in $d,k,r$ are
polylogarithmic in $n$, because $r=O_\eta(k)$ and $k\le K$.
The edge count is therefore
$\Ot_\eta((f+1)^{1+\eta}n^{1+1/k})$, which meets the $\delta$ bound.
Also $d\ge\sqrt{k}/2$ for $k\ge4$, so~\eqref{eq:stretch} has
$r/d=O_\eta(\sqrt{k})$ and $r=O_\eta(k)$.
The construction is deterministic on every input; no expected-time or
random-success qualification is needed. Restoring the original capped
parameters proves all claims.
\end{proof}

\section{Verification and limitations}

Exact finite checks supplement the proof. A reference implementation
constructed representative families, ran the protected scan, and projected
all pendant scales. It tested every labeled graph on $n=2,3,4$ and ten
fixed graph masks each on $n=5,6$, for
$d\le\min\{3,n-1\}$, $f\le\min\{2,n-2\}$ and distinct radii among
$d,d+1,2d+1,48d^2$. All 3,160 mixed constructions passed, including
6,248 individual scans, 23,909 insertion-witness checks, 809,892
rejection-certificate checks and 57,856 final mixed-distance comparisons.
The largest augmented graph had 18 vertices. Separately, 20,465 rational
parameter cases checked the displayed packing inequalities. A separate
implementation checked representative avoidance and distances on 84 small
graphs, with no failures. The proofs establish the asymptotic guarantees,
while the scripts check these finite instances. OpenAI Codex materially
assisted development of the algorithm and proofs and implementation and
checking of these finite verifications.

The algorithm realizes the full mixed-stretch tradeoff in deterministic
polynomial time using the prior tools credited above. The polynomial
exponent is unoptimized and can grow as $\delta$ decreases.

\section*{AI methodology}
OpenAI Codex materially assisted all sections through literature comparison,
algorithm and proof development, implementation, checking and writing.
Separate Codex contexts reviewed the argument and primary sources; the
accompanying scripts perform the finite verification described above.

\bibliographystyle{ACM-Reference-Format}
\bibliography{references}
\end{document}
